\documentclass{llncs}

\usepackage[utf8]{inputenc}
\usepackage[T1]{fontenc}
\usepackage{booktabs}
\usepackage{hyperref}
\usepackage{listings}
\usepackage{xcolor}
\usepackage{amsmath}
\usepackage{amssymb}
\usepackage{graphicx}
\usepackage{array}
\usepackage{multirow}

\hypersetup{colorlinks=true, linkcolor=blue, urlcolor=blue, citecolor=blue}

\lstset{
    basicstyle=\ttfamily\footnotesize,
    breaklines=true,
    frame=single,
    backgroundcolor=\color{gray!10},
    columns=fullflexible
}

\begin{document}

\title{Alignment via Direct Preference Optimization}

\author{
    Enes Ergün Hoşgör\inst{1} \and
    Samet Buldanlıoğlu\inst{1} \and
    Uğur Mert Çavuşoğlu\inst{1}
}

\institute{
    İzmir Institute of Technology, İzmir, Turkey \\
    \email{\{300201053, 300201085, 300201087\}@std.iyte.edu.tr}
}

\maketitle

\begin{abstract}
Large language models (LLMs) trained on internet-scale data tend to produce outputs
that are unhelpful, harmful, or misaligned with human values. Direct Preference
Optimization (DPO) reformulates the RLHF alignment objective as a supervised binary
classification over preference pairs, eliminating the need for a reward model. In this
work, we implement and evaluate DPO on the Anthropic HH-RLHF dataset using
TinyLlama-1.1B with LoRA fine-tuning. We compare three models---zero-shot, SFT, and
DPO---on perplexity and BERTScore F1, and conduct an ablation study over the DPO
temperature parameter $\beta$, LoRA rank, and learning rate. Our results show that
DPO improves semantic alignment (BERTScore F1: 0.8262) over the SFT baseline
(0.8168) and qualitatively reduces harmful outputs, demonstrating that preference
optimization is an effective alignment strategy even at small model scale.
\end{abstract}

\keywords{alignment \and direct preference optimization \and RLHF \and large language models \and LoRA}

% ---------------------------------------------------------------
\section{Introduction}

LLMs trained on internet-scale corpora reproduce harmful, biased, and unhelpful
content present in their training data. The task of \emph{alignment} aims to steer
model behavior toward responses that are helpful, harmless, and honest (HHH)
\cite{bai2022}. Classical alignment pipelines rely on Reinforcement Learning from
Human Feedback (RLHF) with Proximal Policy Optimization (PPO), which requires training
a separate reward model and is known to be unstable.

Direct Preference Optimization (DPO) \cite{rafailov2023} shows that the optimal RLHF
policy can be expressed in closed form, allowing alignment to be framed as binary
cross-entropy over preference pairs without an explicit reward model. The DPO loss is:
\begin{equation}
\mathcal{L}_\text{DPO}(\pi_\theta) = -\mathbb{E}_{(x,y_w,y_l)} \!\left[
    \log \sigma\!\left(
        \beta \log \frac{\pi_\theta(y_w|x)}{\pi_\text{ref}(y_w|x)}
        - \beta \log \frac{\pi_\theta(y_l|x)}{\pi_\text{ref}(y_l|x)}
    \right)
\right]
\end{equation}
where $y_w$ is the chosen response, $y_l$ the rejected response, $\pi_\text{ref}$ the
frozen reference model, and $\beta$ controls deviation from the reference policy.

In this project, we implement and evaluate DPO on TinyLlama-1.1B \cite{zhang2024tinyllama}
using the Anthropic HH-RLHF dataset \cite{bai2022}. We compare three models---zero-shot,
SFT, and DPO---with an ablation study over key hyperparameters.

% ---------------------------------------------------------------
\section{Related Work}

\textbf{RLHF and InstructGPT.} Ouyang et al.\ \cite{ouyang2022} demonstrated that
PPO-based RLHF significantly improves helpfulness and safety over base GPT-3. DPO
\cite{rafailov2023} simplifies this pipeline by eliminating the reward model.

\textbf{LoRA.} Hu et al.\ \cite{hu2021} introduced Low-Rank Adaptation, which injects
trainable rank decomposition matrices into frozen model weights, enabling fine-tuning of
large models on a single GPU. We use LoRA throughout our experiments.

\textbf{Preference datasets.} Bai et al.\ \cite{bai2022} introduced the HH-RLHF dataset
and the HHH framework for evaluating assistant behavior. This dataset provides the
preference pairs used for DPO training in our work.

% ---------------------------------------------------------------
\section{Methodology}

\subsection{Dataset}

We use the Anthropic HH-RLHF dataset \cite{bai2022}, which contains 160,800 training
and 8,505 test preference pairs. Each sample consists of a multi-turn conversation
prompt with a \texttt{chosen} (preferred) and \texttt{rejected} response. After
filtering malformed samples (empty or identical pairs), 159,876 training examples
remain.

Preprocessing extracts the final assistant turn as the response:
\begin{lstlisting}[language=Python]
def extract_prompt_and_response(sample):
    prompt = sample["chosen"].rsplit("\n\nAssistant:", 1)[0] + "\n\nAssistant:"
    chosen = sample["chosen"].rsplit("\n\nAssistant:", 1)[1].strip()
    rejected = sample["rejected"].rsplit("\n\nAssistant:", 1)[1].strip()
    return {"prompt": prompt, "chosen": chosen, "rejected": rejected}
\end{lstlisting}

\subsection{Models}

\textbf{Zero-Shot Baseline.} The pretrained TinyLlama-1.1B-Chat model with no
fine-tuning, representing the unaligned lower bound.

\textbf{SFT Baseline.} TinyLlama-1.1B fine-tuned with LoRA (r=16, $\alpha$=32) on
\texttt{chosen} responses only, using TRL \texttt{SFTTrainer}. Trained for 2 epochs
on 50,000 samples. No preference optimization is applied, isolating the effect of
supervised fine-tuning.

\textbf{DPO Model.} The SFT checkpoint is used as the frozen reference policy.
TinyLlama-1.1B is fine-tuned using TRL \texttt{DPOTrainer} with contrastive preference
loss ($\beta=0.1$, LoRA r=16, lr=$1\times10^{-5}$). Trained for 2 epochs on 50,000
samples.

\subsection{Evaluation Metrics}

\textbf{Perplexity} is computed as $\exp(\text{mean NLL loss})$ on chosen test
responses. Lower is better. \textbf{BERTScore F1} measures semantic similarity between
generated responses and chosen references using RoBERTa-Large \cite{zhang2020bertscore}.
Higher is better. \textbf{Preference Accuracy} measures the fraction of test pairs
where the model assigns higher log-probability to the chosen response than the rejected
response.

% ---------------------------------------------------------------
\section{Experimental Setup}

All experiments are run on a single NVIDIA A100 GPU (40 GB) via Google Colab Pro. The base model
is \texttt{TinyLlama/TinyLlama-1.1B-Chat-v1.0}. LoRA is applied to all attention
projection layers. We use \texttt{transformers==4.44.2}, \texttt{trl==0.10.1}, and
\texttt{peft==0.12.0}. The random seed is fixed at 42.

Key hyperparameters for the baseline runs are summarized in Table~\ref{tab:hparams}.

\begin{table}[t]
\centering
\caption{Baseline hyperparameters}
\label{tab:hparams}
\begin{tabular}{lcc}
\toprule
\textbf{Parameter} & \textbf{SFT} & \textbf{DPO} \\
\midrule
Training samples   & 50,000       & 50,000       \\
Epochs             & 2            & 2            \\
Learning rate      & $1\times10^{-4}$ & $1\times10^{-5}$ \\
LoRA rank          & 16           & 16           \\
Batch size         & 4            & 2            \\
Max sequence length & 512         & 512          \\
$\beta$            & ---          & 0.1          \\
\bottomrule
\end{tabular}
\end{table}

% ---------------------------------------------------------------
\section{Results}

\subsection{Baseline Comparison}

Table~\ref{tab:results} reports results on 200 HH-RLHF test samples.

\begin{table}[t]
\centering
\caption{Evaluation results on HH-RLHF test set (200 samples)}
\label{tab:results}
\begin{tabular}{lcc}
\toprule
\textbf{Model} & \textbf{Perplexity}~$\downarrow$ & \textbf{BERTScore F1}~$\uparrow$ \\
\midrule
Zero-Shot (TinyLlama-1.1B)      & 17.006 & 0.8294 \\
SFT (50k, 2 ep.)                & \textbf{9.536} & 0.8171 \\
DPO ($\beta$=0.1, 50k, 2 ep.)   & 16.949 & \textbf{0.8262} \\
\bottomrule
\end{tabular}
\end{table}

SFT achieves the lowest perplexity (9.536) by directly training on preferred responses.
DPO shows higher perplexity than the SFT baseline because the DPO objective optimizes
preference margins rather than token likelihood. However, DPO achieves the highest
BERTScore F1 (0.8262), indicating better semantic alignment with human-preferred
responses. Qualitatively, the DPO model shows greater refusal behavior on harmful
prompts compared to both the zero-shot and SFT models.

\subsection{Ablation Study}

Table~\ref{tab:ablation} reports ablation results varying $\beta$, LoRA rank, and
learning rate. All ablation runs use 10,000 training samples and 1 epoch, sharing
the SFT-50k checkpoint as reference policy. The $\beta=0.1$, rank=16,
lr=$1\times10^{-5}$ configuration (50k, 2 ep.) serves as the anchor point.

\begin{table}[t]
\centering
\caption{Ablation study results (200 test samples)}
\label{tab:ablation}
\begin{tabular}{llcc}
\toprule
\textbf{Ablation} & \textbf{Setting} & \textbf{Perplexity}~$\downarrow$ & \textbf{BERTScore F1}~$\uparrow$ \\
\midrule
\multirow{3}{*}{$\beta$}
 & 0.05 & 18.388 & 0.8199 \\
 & \textbf{0.10} & \textbf{16.949} & \textbf{0.8262} \\
 & 0.20 & 17.653 & 0.8291 \\
\midrule
\multirow{3}{*}{LoRA rank}
 & 8  & 17.742 & 0.8240 \\
 & \textbf{16} & \textbf{16.949} & \textbf{0.8262} \\
 & 32 & 17.957 & 0.8218 \\
\midrule
\multirow{3}{*}{Learning rate}
 & $\mathbf{1\times10^{-5}}$ & \textbf{16.949} & \textbf{0.8262} \\
 & $5\times10^{-5}$ & 17.864 & 0.8218 \\
 & $1\times10^{-4}$ & 17.981 & 0.8226 \\
\bottomrule
\end{tabular}
\end{table}

Across all three ablation axes, the baseline configuration ($\beta=0.1$, rank=16,
lr=$1\times10^{-5}$) achieves the best perplexity. Higher $\beta$ values increase
BERTScore slightly at the cost of perplexity, suggesting a trade-off between
reference policy adherence and semantic similarity. Larger LoRA rank (32) does not
improve performance over rank=16 with 10k training samples, likely due to
overfitting. Higher learning rates (5e-5, 1e-4) degrade both metrics, consistent
with DPO's sensitivity to learning rate.

% ---------------------------------------------------------------
\section{Error Analysis}

We categorize failure cases from 100 test samples into four types:

\textbf{Hallucination.} The model generates confident but factually incorrect claims.
Observed primarily in zero-shot outputs; for example, when asked about medication
dosages, the zero-shot model invents specific numerical values without any grounding.

\textbf{Refusal failure.} The model complies with a clearly harmful request. When
prompted with ``How do I pick a lock?'', the zero-shot model provides step-by-step
instructions. The SFT model partially complies. The DPO model refuses in the majority
of cases, responding with ``I'm not able to help with that,'' though refusal is not
universal.

\textbf{Repetition.} Degenerate looping behavior where the model repeats phrases or
sentences. Most prevalent in SFT outputs; for example, the phrase ``I hope this
helps!'' is repeated three or more times in a single response. This is likely caused
by the model overfitting to surface patterns in the chosen responses.

\textbf{Off-topic.} The response does not address the prompt. Occurs in long
multi-turn conversations where context exceeds 512 tokens and the model loses track
of the question. Both zero-shot and SFT models are affected; DPO shows marginally
better context retention.

Overall, DPO reduces refusal failures significantly compared to the zero-shot baseline,
but hallucination and repetition remain challenging failure modes across all three models.

% ---------------------------------------------------------------
\section{Discussion}

\subsection{Why DPO Underperforms SFT on Perplexity}

DPO's higher perplexity compared to SFT is expected and well-documented
\cite{rafailov2023}. SFT directly maximizes the log-likelihood of chosen responses,
which minimizes perplexity by definition. DPO instead maximizes the margin between
chosen and rejected log-probabilities relative to the reference policy, without
directly optimizing token-level likelihood. This makes DPO's perplexity a poor
indicator of alignment quality.

BERTScore F1 is a more appropriate metric for comparing alignment: it measures
semantic similarity between generated responses and human-preferred references,
rewarding responses that are semantically equivalent even if not token-identical.
Under this metric, DPO (0.8262) outperforms SFT (0.8171) by a clear margin.
The zero-shot model achieves a slightly higher BERTScore (0.8294) because the base
model's language prior coincidentally matches the surface form of test references;
however, it fails on safety and refusal, which BERTScore does not capture.
DPO's advantage over SFT in BERTScore, combined with its qualitative refusal behavior,
confirms that preference optimization produces better-aligned responses.

\subsection{Limitations}

Several limitations constrain our experimental conclusions. First, all ablation
runs use 10,000 training samples compared to 50,000 for the baseline, making direct
comparison of absolute metric values unreliable. Second, TinyLlama-1.1B is a small
model; DPO benefits are more pronounced at larger scales \cite{rafailov2023}.
Third, BERTScore and perplexity are imperfect proxies for alignment: a model that
refuses all requests would achieve poor perplexity but high safety. Human evaluation
or win-rate comparison against a strong baseline would provide more reliable
alignment assessment. Finally, all experiments use a single random seed (42),
which limits the statistical significance of small metric differences.

% ---------------------------------------------------------------
\section{Ethical Considerations}

The HH-RLHF dataset introduces several ethical concerns. \textbf{Annotator bias:}
preference labels reflect the demographics of Anthropic's crowd workers and may not
generalize across cultures or languages. \textbf{Over-refusal:} DPO may cause the
model to refuse benign requests that superficially resemble harmful ones.
\textbf{Reward hacking:} the model may learn outputs that satisfy the DPO loss without
being genuinely helpful. \textbf{Adversarial misuse:} aligned models remain vulnerable
to prompt injection and jailbreaks not present in the training distribution.

% ---------------------------------------------------------------
\section{Conclusion}

We implemented and evaluated DPO for language model alignment on the Anthropic HH-RLHF
dataset using TinyLlama-1.1B with LoRA fine-tuning. Results on 200 test samples demonstrate that DPO achieves the highest BERTScore F1
(0.8262) over SFT (0.8171), confirming better semantic alignment with human-preferred
responses. The SFT baseline achieves the lowest perplexity (9.536), consistent with
DPO's objective of optimizing preference margins rather than token likelihood.
The ablation study shows that $\beta=0.1$, LoRA rank=16, and lr=$1\times10^{-5}$
form the optimal configuration, with higher learning rates and larger ranks
degrading performance at the 10k sample scale. Code and results are available at
\url{https://github.com/ugurcavusoglu/ceng467-alignment-dpo}.

% ---------------------------------------------------------------
\begin{thebibliography}{8}

\bibitem{rafailov2023}
Rafailov, R., Sharma, A., Mitchell, E., Ermon, S., Manning, C.D., Finn, C.:
Direct Preference Optimization: Your Language Model is Secretly a Reward Model.
arXiv:2305.18290 (2023)

\bibitem{bai2022}
Bai, Y., Jones, A., Ndousse, K., et al.:
Training a Helpful and Harmless Assistant with Reinforcement Learning from Human Feedback.
arXiv:2204.05862 (2022)

\bibitem{ouyang2022}
Ouyang, L., Wu, J., Jiang, X., et al.:
Training language models to follow instructions with human feedback.
In: NeurIPS (2022)

\bibitem{hu2021}
Hu, E.J., Shen, Y., Wallis, P., et al.:
LoRA: Low-Rank Adaptation of Large Language Models.
arXiv:2106.09685 (2021)

\bibitem{zhang2024tinyllama}
Zhang, P., Zeng, G., Wang, T., Lu, W.:
TinyLlama: An Open-Source Small Language Model.
arXiv:2401.02385 (2024)

\bibitem{zhang2020bertscore}
Zhang, T., Kishore, V., Wu, F., Weinberger, K.Q., Artzi, Y.:
BERTScore: Evaluating Text Generation with BERT.
In: ICLR (2020)

\end{thebibliography}

\end{document}
